% === Begin Label Data ===
% ID: 159
% 难度星级: 
% 标签: 
% 备注: 
% 组卷引用次数: 0
% === End  Label Data ===

\begin{problem}{2023}{G}{上海卷}{21}{数列}

已知 $f(x)=\ln x$, 在该函数图像 $\Gamma$ 上取一点 $a_1$, 过点 $(a_1, f(a_1))$ 做函数 $f(x)$ 的切线, 该切线与 $y$ 轴的交点记作 $(0, a_2)$, 若 $a_2>0$, 则过点 $(a_2, f(a_2))$ 做函数 $f(x)$ 的切线, 该切线与 $y$ 轴的交点记作 $(0, a_3)$, 以此类推 $a_3, a_4, \cdots$, 直至 $a_m \le 0$ 停止, 由这些项构成数列 $\{a_n\}$.

(1) 若正整数 $m>2$，证明 $a_m=\ln a_{m-1}-1$;

(2) 若正整数 $m \geqslant 2$，试比较 $a_m$ 与 $a_{m-1}-2$ 大小;

(3) 若正整数 $k \geqslant 3$，是否存在 $k$ 使得 $a_1, a_2, \cdots, a_k$ 依次成等差数列？若存在，求出 $k$ 的所有取值，若不存在，试说明理由.

\end{problem}

\begin{solutions}
（1）$f'(x) = \dfrac{1}{x}$, 则 $f(x)$ 在 $(a_n, f(a_n))$ 处的切线为 $y - \ln a_n = \dfrac{1}{a_n} (x - a_n)$,

当 $x=0$ 时, $y = \ln a_n - 1$, 即 $a_{n+1} = \ln a_n - 1$,

所以当正整数 $m \ge 2$ 时, $a_m = \ln a_{m-1} - 1$;

（2）作差得 $a_{n+1} - (a_n - 2) = \ln a_n - a_n + 1$,

令 $g(x) = \ln x - x + 1$, $g'(x) = \dfrac{1}{x} - 1$,

当 $0 < x < 1$ 时, $g'(x) > 0$, 当 $x > 1$ 时, $g'(x) < 0$,

故 $g(x)$ 在 $(0, 1)$ 单调递增, 在 $(1, +\infty)$ 上单调递减,

$g(x) \le g(1) = 0$, 故 $a_{n+1} \le a_n - 2$,

所以当正整数 $m \ge 2$ 时, $a_m \le a_{m-1} - 2$;

（3）存在. 理由如下:

假设存在符合要求的 $k$. 由 (2) 可知, 数列 $\{a_n\}$ 为严格的递减数列, $n=1,2,3,\cdots,k$.

由 (1) 可知, 公差 $d = a_n - a_{n-1} = \ln a_{n-1} - a_{n-1} - 1$, $2 \leqslant n \leqslant k$.

令 $h(x) = \ln x - x - 1$, 则 $h'(x) = \dfrac{1}{x} - 1$. 当 $0 < x < 1$ 时, $h'(x) > 0$, $h(x)$ 单调递增; 当 $x > 1$ 时, $h'(x) < 0$, $h(x)$ 单调递减.

则 $h(x) = d$ 至多只有两个解, 即至多存在两个 $a_{n-1}$, 使得 $h(a_{n-1}) = d$.

若 $k \geqslant 4$, 则 $h(a_1) = h(a_2) = h(a_3) = d$, 矛盾, 则 $k = 3$.

当 $k = 3$ 时, 若 $a_1, a_2, a_3$ 依次成等差数列, 则 $2a_2 = a_1 + a_3$.

因为 $a_2 = \ln a_1 - 1, a_3 = \ln a_2 - 1 = \ln(\ln a_1 - 1) - 1$,

所以 $2(\ln a_1 - 1) = a_1 + \ln(\ln a_1 - 1) - 1$,

即 $\ln(\ln a_1 - 1) - 2\ln a_1 + a_1 + 1 = 0$.

设函数 $h(x) = \ln(\ln x - 1) - 2\ln x + x + 1$.

因为 $h(e^{1.1}) = \ln 0.1 - 2.2 + e^{1.1} + 1 = e^{1.1} - \ln 10 - 1.2 < 0, h(e^2) = -3 + e^2 > 0$,

所以存在 $x_0 \in (e^{1.1}, e^2)$, 使得 $h(x_0) = 0$,

所以取 $a_1 = x_0, a_2 = \ln a_1 - 1, a_3 = \ln a_2 - 1$, 它们依次构成等差数列.

综上, $k = 3$.

法二：假设存在正整数 $k > 3$, 使得 $a_1, a_2, a_3, \cdots, a_k$ 依次成等差数列, 公差为 $d$.

则 $a_k - a_{k-1} = \ln a_{k-1} - \ln a_{k-2} = d$, 所以 $\dfrac{a_{k-1}}{a_{k-2}} = e^d$, 所以数列 $\{a_n\}$ 既为等差数列又为等比数列.

故此时 $d = 0$, 与 $a_k - a_{k-1} \leqslant -2$ 矛盾, 所以不存在正整数 $k > 3$, 使得 $a_1, a_2, a_3, \cdots, a_k$ 依次成等差数列.

当 $k = 3$ 时, $a_2 = \ln a_1 - 1$, 所以 $a_1 = e^{a_2+1}, a_3 = \ln a_2 - 1$. 若满足 $a_1 + a_3 = 2a_2$, 则 $e^{a_2+1} + \ln a_2 - 1 - 2a_2 = 0$.

令 $q(x) = e^{x+1} + \ln x - 1 - 2x \ (x > 0)$, 则 $q(1) = e^2 - 3 > 0, q(e^{-10}) = e^{-10+1} - 10 - 1 - 2e^{-10} < e^2 - 10 - 1 < 0$, 所以存在 $x_0 \in (e^{-10}, 1)$, 使得 $q(x_0) = 0$, 即 $a_1 + a_3 = 2a_2$, 所以存在 $k = 3$, 使得 $a_1, a_2, a_3, \cdots, a_k$ 依次成等差数列.

综上, $k = 3$.

故存在最多 3 项成等差数列.
\end{solutions}